% LGS: Lead-biased Grounding Score — a reference-free, embedding-only
% summary-quality metric.
%
% Target venue: IEEE Access.
% This file uses the bare IEEEtran class so it compiles on a stock
% TeX Live install. For the camera-ready submission, swap the
% documentclass to the IEEE Access template (`\documentclass{ieeeaccess}`
% with the IEEEaccess.cls bundled in the official Overleaf template) and
% paste the body below into that template's main file. The body is
% deliberately written to be portable between the two.

\documentclass[journal]{IEEEtran}

\usepackage[utf8]{inputenc}
\usepackage[T1]{fontenc}
\usepackage{amsmath,amssymb}
\usepackage{graphicx}
\usepackage{booktabs}
\usepackage{multirow}
\usepackage{url}
\usepackage{cite}
\usepackage{hyperref}

\hyphenation{summ-eval embed-ding-only ref-erence-free}

\begin{document}

\title{LGS: A Reference-Free, Embedding-Only Metric for Summary
Quality with a Lead-Biased Grounding Prior}

\author{%
  Mikolaj Semeniuk%
  \thanks{Manuscript prepared for IEEE Access. Corresponding author:
  Mikolaj Semeniuk (mikolaj.adam.semeniuk@gmail.com).}%
}

\markboth{IEEE Access, Vol.~XX, No.~X, May~2026}%
{Semeniuk: LGS --- Lead-Biased Grounding Score}

\maketitle

% ------------------------------------------------------------------
\begin{abstract}
We propose \textbf{LGS} (Lead-biased Grounding Score), a
reference-free, embedding-only metric for summary quality. Given a
source article and a candidate summary, LGS is computed as
$\mathrm{score} = \mathrm{mean}_j\, \max_i\, w(i)\cdot
\cos(c_j, s_i)$, where $c_j$ and $s_i$ are sentence embeddings of
the candidate and source respectively, and
$w(i)=\exp(-\lambda\, i/n)$ is an exponential lead-bias prior on
source-sentence position. The single hyperparameter $\lambda$ is
selected on a held-out development split of \textsc{SummEval} and
verified on a disjoint test split. Unlike reference-based metrics
(BERTScore, BARTScore, MoverScore, ROUGE-L, ChrF), LGS does not
require a human reference summary, and unlike the reference-free
G-Eval it does not require a large language-model judge call per
candidate. The metric runs with one sentence-embedder forward pass
per source/candidate sentence (we use \texttt{nomic-embed-text}, a
137M-parameter embedder served locally through Ollama). On
\textsc{SummEval}, LGS significantly outperforms a same-embedder
whole-text cosine baseline on coherence, consistency, and fluency
($\Delta\rho=+.117/+.171/+.116$, $p<.001$), and is competitive with
or stronger than reference-based embedding metrics
(BERTScore, SMART-Model) on multiple dimensions. We do not claim
state of the art: UniEval and G-Eval remain stronger overall, and
our headline contribution is the combination of a reference-free
formulation, a single-embedder cost profile, and a held-out
methodology that we show is necessary in practice because the
optimal $\lambda$ shifts across embedder backbones.
\end{abstract}

\begin{IEEEkeywords}
Summarization evaluation, reference-free evaluation, sentence
embeddings, SummEval, lead bias, automatic metrics.
\end{IEEEkeywords}

% ------------------------------------------------------------------
\section{Introduction}
\label{sec:intro}

% TODO(prose): tighten, motivate practical deployment context.
Automatic evaluation of generated summaries remains an open problem
even after a decade of progress on reference-based n-gram and
embedding metrics. The two strongest published metrics on the
\textsc{SummEval} benchmark~\cite{fabbri2021summeval} --- UniEval and
G-Eval --- both depend on substantial inference cost: UniEval
requires a fine-tuned T5-large pass per (article, candidate, dimension)
tuple, and G-Eval queries a large language-model judge with a
chain-of-thought prompt and aggregates over multiple samples. In
deployment scenarios where evaluation must run online over many
candidates (model selection, RLHF reward modelling, regression
testing of summarisation services), this cost matters.

We propose a metric, \textbf{LGS} (Lead-biased Grounding Score),
that fills a different niche: \emph{reference-free} (it requires
only the source article and the candidate, not a human reference),
\emph{embedding-only} (one forward pass per sentence with a small
sentence embedder; no LLM judge, no fine-tuned QA model), and
\emph{interpretable} (each candidate sentence is anchored to its
single best-matching source sentence under a position prior). The
contribution of this paper is not a new state of the art on
\textsc{SummEval} --- UniEval and G-Eval remain stronger on the
dimensions where each excels --- but a metric that
\begin{enumerate}
\item is reference-free, removing a hard prerequisite (human
reference summaries) that most published embedding metrics share;
\item runs at a fraction of the cost of LLM-judge or
fine-tuned-encoder metrics;
\item is selected with a held-out development/test methodology that,
as we demonstrate empirically, is \emph{necessary} for this metric
because the optimal lead-bias exponent shifts non-trivially across
sentence-embedder backbones.
\end{enumerate}

The rest of the paper is structured as follows.
Section~\ref{sec:related} surveys reference-based and reference-free
metrics and positions LGS against the closest neighbours in each
family. Section~\ref{sec:method} defines LGS formally and motivates
each design choice. Section~\ref{sec:experiments} describes the
experimental protocol. Section~\ref{sec:results} reports
correlations with human judgment on \textsc{SummEval} together
with paired-bootstrap significance tests against thirteen baselines.
Section~\ref{sec:ablation} reports the lead-bias hyperparameter
sweep and the embedder ablation. Section~\ref{sec:discussion}
discusses limitations and the cost trade-off vs G-Eval and UniEval.
Section~\ref{sec:conclusion} concludes.

% ------------------------------------------------------------------
\section{Related Work}
\label{sec:related}

% TODO(prose): expand each subsection. Currently a structural skeleton.

\subsection{Reference-based metrics}
N-gram metrics (BLEU~\cite{papineni2002bleu},
ROUGE-L~\cite{lin2004rouge}, METEOR~\cite{banerjee2005meteor},
ChrF~\cite{popovic2015chrf}) and contextual-embedding metrics
(BERTScore~\cite{zhang2020bertscore},
MoverScore~\cite{zhao2019moverscore},
BARTScore~\cite{yuan2021bartscore},
SMART~\cite{amplayo2022smart}) all require one or more human
reference summaries. The dependence is structural: the score is a
function of \texttt{(reference, candidate)} pairs, not of the source
article. In production deployments, reference summaries rarely
exist; this is the primary gap LGS addresses.

\subsection{Reference-free / source-based metrics}
G-Eval~\cite{liu2023geval} prompts a large language model to score
each candidate against the source. It is reference-free and obtains
the strongest correlations on \textsc{SummEval}'s consistency
dimension among published metrics, but the cost profile (one or
more LLM-judge calls per candidate) and the stochasticity of
sampled decoding make it unsuitable for scenarios that require
large-scale, deterministic scoring.

UniEval~\cite{zhong2022unieval} is a hybrid: it runs a fine-tuned
T5-large encoder over a Boolean question (``Is this summary
\{coherent,consistent,fluent,relevant\}?'') and computes the
posterior of ``Yes''. This requires a fine-tuned model per
dimension and a forward pass per (article, candidate, dimension).

\subsection{The closest baseline: whole-text embedding cosine}
The simplest reference-free embedding baseline is a single cosine
similarity between the embedding of the entire source and the
embedding of the entire candidate. We call this \textsc{EmbedScorer}
and use it as the headline like-for-like baseline: it shares the
embedder backbone with LGS, so any improvement is directly
attributable to the sentence-level grounding plus lead-bias prior
contribution rather than to embedder choice.

% ------------------------------------------------------------------
\section{Method}
\label{sec:method}

\subsection{Definition}
Let $D$ be a source document split into $n$ sentences
$s_0,\dots,s_{n-1}$ and let $C$ be a candidate summary split into
$m$ sentences $c_0,\dots,c_{m-1}$. Let $\mathrm{emb}(\cdot)$ be a
sentence embedder (we use \texttt{nomic-embed-text} served via
Ollama, see Section~\ref{sec:experiments}). The Lead-biased
Grounding Score is
\begin{equation}
\label{eq:lgs}
\mathrm{LGS}(D,C) \;=\; \frac{1}{m}\sum_{j=0}^{m-1}
\;\max_{0\le i<n}\; w(i)\cdot \cos\!\bigl(\mathrm{emb}(c_j),
\mathrm{emb}(s_i)\bigr),
\end{equation}
where the source-position prior is
\begin{equation}
\label{eq:wprior}
w(i) \;=\; \exp\!\bigl(-\lambda \cdot i / n\bigr),
\end{equation}
$i=0$ denotes the lead sentence, and $\lambda \ge 0$ is the only
hyperparameter. $\lambda=0$ disables the prior and reduces LGS to
position-agnostic mean-of-max recall.

\subsection{Design rationale}
% TODO(prose): expand each bullet into a paragraph.
\begin{itemize}
\item \textbf{Sentence-level grounding}: a candidate sentence is
``grounded'' if some source sentence supports it. Aggregating with
$\max_i$ on the source side and $\mathrm{mean}_j$ on the candidate
side gives a localised hallucination-sensitive signal: a single
ungrounded candidate sentence depresses the score by $1/m$.

\item \textbf{Why exponential lead-bias?} News summarisation
corpora (notably CNN/DailyMail, the source of all
\textsc{SummEval} articles) exhibit well-documented lead bias:
salient information is concentrated in the first few sentences.
The exponential family $\exp(-\lambda i/n)$ is the simplest decay
shape that interpolates between uniform ($\lambda=0$) and a
near-step prior ($\lambda \to \infty$); the prior is normalised by
document length so the same $\lambda$ behaves consistently across
short and long articles.

\item \textbf{Why a single hyperparameter?} Earlier formulations of
this metric included additional terms --- a precision side, a
coverage exponent, a within-summary diversity penalty, top-$k$
recall --- each of which was tested on the held-out development
split and rejected for failing to lift the dev mean Spearman
$\rho$. A development log of these rejected components is preserved
in the project repository for reproducibility but is not part of
the canonical metric.
\end{itemize}

\subsection{Held-out hyperparameter selection}
\label{sec:method-heldout}
\textsc{SummEval} contains 100 source articles, each summarised by
16 systems, for 1{,}600 (source, candidate) pairs with human
ratings on four dimensions (coherence, consistency, fluency,
relevance). We split at the article level, deterministically by
JSONL order: the first 50 articles form the development split, the
remaining 50 the test split. We sweep
$\lambda \in \{0, 0.25, 0.5, 1.0, 2.0\}$ on the development split,
score each variant with summary-level Spearman $\rho$ averaged over
the four dimensions, and select the dev-best $\lambda^*$. The
selected $\lambda^*$ is then re-evaluated on the test split as a
sanity check; the test split is \emph{not} consulted during
selection. The full headline run is computed on all 100 articles
with $\lambda = \lambda^*$.

% ------------------------------------------------------------------
\section{Experimental Setup}
\label{sec:experiments}

% TODO(prose): expand to full subsections.
\subsection{Dataset}
\textsc{SummEval}~\cite{fabbri2021summeval}: 100 CNN/DailyMail
articles $\times$ 16 abstractive summarisation systems
$=$ 1{,}600 (source, candidate) pairs. Each pair has expert human
ratings on coherence (Coh), consistency (Con), fluency (Flu) and
relevance (Rel) on a 1--5 scale. We use summary-level correlations
as the primary metric (Spearman $\rho$, Kendall $\tau$); system-level
correlations are reported as a secondary check.

\subsection{Embedder}
The canonical run of LGS uses \texttt{nomic-embed-text}
(137M parameters), served locally via Ollama. We additionally run
the embedder ablation in Section~\ref{sec:ablation} with
\texttt{all-minilm} (23M), \texttt{mxbai-embed-large} (335M), and
\texttt{bge-m3} (567M) to verify that the metric design is not
specific to one backbone.

\subsection{Baselines}
We compare against thirteen baselines spanning the reference-based,
reference-free, and mixed families:
reference-based: BLEU, ROUGE-L, ChrF, METEOR, SMART-String,
SMART-Model, EmbedScorer (whole-text cosine, ours, same embedder
as LGS for like-for-like), BERTScore, MoverScore, BARTScore,
GPTScore;
source-based: G-Eval; mixed: UniEval. All baselines use simplified
implementations relative to the original papers (uniform IDF for
MoverScore; GPT-2 instead of GPT-3 for GPTScore; greedy
single-sample decoding for G-Eval; single ``Yes/No'' query for
UniEval). These simplifications are documented per metric and do
not change the ordering of methods qualitatively.

\subsection{Significance testing}
For each (LGS, baseline, dimension) cell we run a paired bootstrap
over articles (10{,}000 resamples) on summary-level Spearman, and
report the mean delta, the 95\% confidence interval, and a two-sided
$p$-value.

% ------------------------------------------------------------------
\section{Results}
\label{sec:results}

% TODO(prose): expand discussion under each table.
Table~\ref{tab:correlations_summary} reports summary-level Spearman
$\rho$ and Kendall $\tau$ correlations with human judgment for all
fourteen metrics. Table~\ref{tab:compare_lgs_summary} reports the
paired-bootstrap deltas of LGS against each baseline on Spearman
$\rho$. Table~\ref{tab:correlations_system} reports system-level
correlations (note that with $n=16$ systems the confidence intervals
are wide; we treat this table as a secondary sanity check).

\begin{table*}[t]
\centering
\caption{Summary-level correlations with human judgment on SummEval. $\rho$ = Spearman, $r$ = Pearson, $\tau$ = Kendall.}
\label{tab:correlations_summary}
\begin{tabular}{lrrrrrrrrrrrr}
\toprule
Metric & \multicolumn{3}{c}{Coh} & \multicolumn{3}{c}{Con} & \multicolumn{3}{c}{Flu} & \multicolumn{3}{c}{Rel} \\
\cmidrule(lr){2-4} \cmidrule(lr){5-7} \cmidrule(lr){8-10} \cmidrule(lr){11-13} 
  & $\rho$ & $r$ & $\tau$ & $\rho$ & $r$ & $\tau$ & $\rho$ & $r$ & $\tau$ & $\rho$ & $r$ & $\tau$ \\
\midrule
BLEU & .102 {\scriptsize [.036, .165]} & .107 {\scriptsize [.045, .161]} & .072 {\scriptsize [.024, .115]} & .137 {\scriptsize [.085, .185]} & .151 {\scriptsize [.109, .195]} & .107 {\scriptsize [.065, .147]} & .066 {\scriptsize [.002, .123]} & .106 {\scriptsize [.057, .158]} & .051 {\scriptsize [.005, .097]} & .206 {\scriptsize [.123, .277]} & .204 {\scriptsize [.123, .288]} & .148 {\scriptsize [.094, .203]} \\
ROUGE-L & .163 {\scriptsize [.093, .231]} & .166 {\scriptsize [.093, .227]} & .115 {\scriptsize [.067, .162]} & .140 {\scriptsize [.085, .195]} & .159 {\scriptsize [.111, .203]} & .110 {\scriptsize [.062, .154]} & .108 {\scriptsize [.043, .167]} & .146 {\scriptsize [.090, .206]} & .083 {\scriptsize [.036, .134]} & .201 {\scriptsize [.115, .278]} & .210 {\scriptsize [.132, .287]} & .144 {\scriptsize [.084, .202]} \\
ChrF & .158 {\scriptsize [.092, .219]} & .137 {\scriptsize [.070, .204]} & .112 {\scriptsize [.065, .160]} & .174 {\scriptsize [.125, .220]} & .199 {\scriptsize [.149, .252]} & .137 {\scriptsize [.101, .174]} & .110 {\scriptsize [.050, .167]} & .149 {\scriptsize [.087, .211]} & .086 {\scriptsize [.043, .132]} & .293 {\scriptsize [.219, .361]} & .330 {\scriptsize [.256, .403]} & .210 {\scriptsize [.158, .264]} \\
METEOR & .199 {\scriptsize [.146, .250]} & .206 {\scriptsize [.153, .259]} & .139 {\scriptsize [.100, .176]} & .134 {\scriptsize [.085, .179]} & .169 {\scriptsize [.125, .213]} & .105 {\scriptsize [.064, .143]} & .130 {\scriptsize [.066, .193]} & .187 {\scriptsize [.133, .246]} & .101 {\scriptsize [.051, .151]} & .185 {\scriptsize [.115, .251]} & .198 {\scriptsize [.126, .267]} & .131 {\scriptsize [.078, .179]} \\
SMART-String & -.006 {\scriptsize [-.070, .063]} & .020 {\scriptsize [-.047, .084]} & -.004 {\scriptsize [-.051, .043]} & .111 {\scriptsize [.054, .166]} & .162 {\scriptsize [.114, .211]} & .087 {\scriptsize [.043, .130]} & .105 {\scriptsize [.046, .159]} & .169 {\scriptsize [.115, .226]} & .081 {\scriptsize [.037, .123]} & .124 {\scriptsize [.045, .194]} & .153 {\scriptsize [.078, .228]} & .088 {\scriptsize [.034, .143]} \\
EmbedScorer & .274 {\scriptsize [.146, .380]} & .311 {\scriptsize [.167, .434]} & .193 {\scriptsize [.104, .282]} & .168 {\scriptsize [.025, .310]} & .132 {\scriptsize [-.030, .280]} & .131 {\scriptsize [.009, .244]} & .150 {\scriptsize [-.083, .351]} & .156 {\scriptsize [-.103, .379]} & .111 {\scriptsize [-.061, .269]} & .247 {\scriptsize [.095, .403]} & .350 {\scriptsize [.125, .528]} & .178 {\scriptsize [.054, .299]} \\
SMART-Model & .357 {\scriptsize [.244, .454]} & .387 {\scriptsize [.264, .488]} & .259 {\scriptsize [.177, .340]} & .227 {\scriptsize [.092, .358]} & .146 {\scriptsize [.003, .281]} & .178 {\scriptsize [.074, .276]} & .155 {\scriptsize [-.035, .331]} & .162 {\scriptsize [-.036, .353]} & .117 {\scriptsize [-.028, .250]} & .325 {\scriptsize [.182, .457]} & .405 {\scriptsize [.213, .556]} & .240 {\scriptsize [.134, .335]} \\
G-Eval & .461 {\scriptsize [.351, .570]} & .476 {\scriptsize [.374, .576]} & .379 {\scriptsize [.289, .474]} & .485 {\scriptsize [.404, .567]} & .574 {\scriptsize [.467, .666]} & .446 {\scriptsize [.372, .523]} & .143 {\scriptsize [.062, .206]} & .129 {\scriptsize [.077, .184]} & .133 {\scriptsize [.062, .191]} & .547 {\scriptsize [.415, .659]} & .624 {\scriptsize [.488, .724]} & .470 {\scriptsize [.351, .571]} \\
\bottomrule
\end{tabular}
\end{table*}

\begin{table*}[t]
\centering
\caption{Paired bootstrap comparison of BGS against baselines on summary-level Spearman correlations. Each cell reports $\Delta\rho$ (ours minus baseline) with 95\% CI in brackets and $p$-value. Bold marks significant improvements ($p<0.05$, CI excludes 0).}
\label{tab:compare_bgs_summary}
\begin{tabular}{lcccc}
\toprule
Baseline & Coh & Con & Flu & Rel \\
\midrule
UniEval & -.173 {\scriptsize [-.236, -.113], $p$=<.001} & -.177 {\scriptsize [-.233, -.123], $p$=<.001} & -.287 {\scriptsize [-.355, -.226], $p$=<.001} & -.107 {\scriptsize [-.176, -.039], $p$=0.003} \\
BERTScore & +.019 {\scriptsize [-.039, +.076], $p$=0.509} & \textbf{+.118 {\scriptsize [+.059, +.181], $p$=<.001}} & +.015 {\scriptsize [-.039, +.068], $p$=0.612} & -.016 {\scriptsize [-.081, +.052], $p$=0.634} \\
G-Eval & +.022 {\scriptsize [-.058, +.101], $p$=0.589} & -.261 {\scriptsize [-.321, -.202], $p$=<.001} & +.040 {\scriptsize [-.032, +.111], $p$=0.273} & -.063 {\scriptsize [-.144, +.020], $p$=0.137} \\
SMART-Model & \textbf{+.068 {\scriptsize [+.020, +.118], $p$=0.005}} & \textbf{+.134 {\scriptsize [+.084, +.182], $p$=<.001}} & \textbf{+.089 {\scriptsize [+.038, +.141], $p$=<.001}} & +.035 {\scriptsize [-.022, +.095], $p$=0.248} \\
ChrF & \textbf{+.237 {\scriptsize [+.145, +.326], $p$=<.001}} & +.053 {\scriptsize [-.013, +.121], $p$=0.123} & +.046 {\scriptsize [-.017, +.113], $p$=0.171} & +.062 {\scriptsize [-.031, +.153], $p$=0.188} \\
\bottomrule
\end{tabular}
\end{table*}

\begin{table*}[t]
\centering
\caption{System-level correlations with human judgment on SummEval. $\rho$ = Spearman, $r$ = Pearson, $\tau$ = Kendall.}
\label{tab:correlations_system}
\begin{tabular}{lrrrrrrrrrrrr}
\toprule
Metric & \multicolumn{3}{c}{Coh} & \multicolumn{3}{c}{Con} & \multicolumn{3}{c}{Flu} & \multicolumn{3}{c}{Rel} \\
\cmidrule(lr){2-4} \cmidrule(lr){5-7} \cmidrule(lr){8-10} \cmidrule(lr){11-13} 
  & $\rho$ & $r$ & $\tau$ & $\rho$ & $r$ & $\tau$ & $\rho$ & $r$ & $\tau$ & $\rho$ & $r$ & $\tau$ \\
\midrule
BLEU & .276 {\scriptsize [-.361, .720]} & .210 {\scriptsize [-.257, .759]} & .183 {\scriptsize [-.276, .607]} & .024 {\scriptsize [-.589, .543]} & .605 {\scriptsize [-.245, .877]} & .017 {\scriptsize [-.407, .459]} & .415 {\scriptsize [-.223, .818]} & .583 {\scriptsize [-.067, .889]} & .283 {\scriptsize [-.228, .691]} & .506 {\scriptsize [-.134, .844]} & .573 {\scriptsize [.071, .830]} & .350 {\scriptsize [-.146, .709]} \\
ROUGE-L & .312 {\scriptsize [-.324, .734]} & .355 {\scriptsize [-.131, .768]} & .217 {\scriptsize [-.243, .636]} & -.018 {\scriptsize [-.607, .526]} & .528 {\scriptsize [-.178, .783]} & -.017 {\scriptsize [-.477, .424]} & .365 {\scriptsize [-.243, .771]} & .552 {\scriptsize [.053, .802]} & .250 {\scriptsize [-.183, .648]} & .482 {\scriptsize [-.139, .851]} & .586 {\scriptsize [.093, .853]} & .317 {\scriptsize [-.143, .712]} \\
ChrF & .526 {\scriptsize [.022, .831]} & .162 {\scriptsize [-.169, .814]} & .350 {\scriptsize [-.045, .673]} & .794 {\scriptsize [.505, .944]} & .717 {\scriptsize [.374, .984]} & .617 {\scriptsize [.351, .852]} & .832 {\scriptsize [.555, .940]} & .638 {\scriptsize [.394, .969]} & .617 {\scriptsize [.357, .835]} & .803 {\scriptsize [.498, .917]} & .622 {\scriptsize [.536, .878]} & .617 {\scriptsize [.375, .804]} \\
METEOR & .094 {\scriptsize [-.547, .582]} & .206 {\scriptsize [-.361, .642]} & .000 {\scriptsize [-.472, .404]} & -.318 {\scriptsize [-.728, .229]} & .195 {\scriptsize [-.420, .543]} & -.200 {\scriptsize [-.568, .264]} & .041 {\scriptsize [-.538, .631]} & .247 {\scriptsize [-.275, .592]} & .033 {\scriptsize [-.439, .523]} & .179 {\scriptsize [-.487, .641]} & .292 {\scriptsize [-.346, .743]} & .067 {\scriptsize [-.439, .486]} \\
SMART-String & -.029 {\scriptsize [-.777, .511]} & -.063 {\scriptsize [-.538, .463]} & -.150 {\scriptsize [-.655, .297]} & -.197 {\scriptsize [-.656, .330]} & .449 {\scriptsize [-.343, .862]} & -.117 {\scriptsize [-.469, .268]} & .171 {\scriptsize [-.402, .679]} & .379 {\scriptsize [-.328, .852]} & .117 {\scriptsize [-.299, .535]} & .256 {\scriptsize [-.468, .735]} & .375 {\scriptsize [-.306, .750]} & .150 {\scriptsize [-.386, .593]} \\
\bottomrule
\end{tabular}
\end{table*}


\subsection{Headline findings}
% TODO(prose): turn into proper paragraphs.
\begin{itemize}
\item LGS \emph{significantly} outperforms the like-for-like
\textsc{EmbedScorer} baseline (whole-text cosine on the same
embedder) on Coh, Con and Flu
($\Delta\rho = +.117 / +.171 / +.116$, all $p<.001$), with a
statistical tie on Rel. This is the cleanest support for the
``sentence-level grounding plus lead-bias prior'' contribution of
the metric.
\item LGS significantly outperforms SMART-Model on Coh, Con and
Flu, BERTScore on Con, and ChrF on Coh. It loses to UniEval on
every dimension and to G-Eval on Con. We treat this as the
honest ceiling.
\end{itemize}

% ------------------------------------------------------------------
\section{Ablation Studies}
\label{sec:ablation}

\subsection{Lead-bias exponent $\lambda$}
\label{sec:ablation-lambda}
% TODO(prose): expand interpretation.
Table~\ref{tab:ablation_lgs} reports the development-split sweep of
$\lambda \in \{0, 0.25, 0.5, 1.0, 2.0\}$ and the test-split
verification of the dev-selected $\lambda^* = 0.5$. The dev mean
Spearman $\rho$ across the four dimensions is .233 at $\lambda=0$
(no prior) and .252 at $\lambda^*=0.5$, a lift of $+.019$. On the
held-out test split, $\lambda^*=0.5$ also outperforms the no-prior
baseline on the mean (.319 vs .314). Beyond $\lambda > 1$ the
dev mean drops, indicating that the prior becomes too aggressive
once the article tail is heavily suppressed.

\begin{table*}[t]
\centering
\caption{Ablation of BGS on SummEval (summary-level Spearman $\rho$). β controls the F$_\beta$ recall vs.\ precision weighting; salience-frac is the top-$k$\% of source sentences (by degree centrality) used in the precision side. Canonical configuration ($\beta=2$, frac $=0.30$) marked $\star$.}
\label{tab:ablation_bgs}
\begin{tabular}{lrrrr}
\toprule
Variant & Coh & Con & Flu & Rel \\
\midrule
$\beta=1$, frac=0.10 & .332 & .149 & .100 & .363 \\
$\beta=1$, frac=0.30 & .341 & .155 & .099 & .359 \\
$\beta=1$, frac=0.50 & .348 & .161 & .108 & .361 \\
$\beta=1$, frac=1.00 & .353 & .161 & .113 & .361 \\
$\beta=2$, frac=0.30 $\star$ & .373 & .203 & .134 & .356 \\
$\beta=3$, frac=0.30 & .377 & .214 & .143 & .348 \\
\midrule
Recall only (no precision side) & .377 & .223 & .149 & .337 \\
\bottomrule
\end{tabular}
\end{table*}


% TODO(reviewer-pre-empt): add finer-grained sweep around the
% optimum (lambda in {0.3, 0.4, 0.5, 0.6, 0.75}) with bootstrap CI
% on each dev-mean rho. The current dev-mean gap between
% lambda=0.25 (.251) and lambda=0.5 (.252) is .001, which is well
% inside bootstrap noise; a finer sweep with CIs forestalls the
% obvious reviewer pushback that lambda*=0.5 is a noise pick.

\subsection{Embedder ablation}
\label{sec:ablation-embedder}
% TODO(prose): expand interpretation, especially the "headline
% nomic is not the highest-mean embedder" honesty paragraph.
Table~\ref{tab:ablation_embedder} runs the canonical $\lambda=0.5$
metric with four sentence-embedder backbones spanning a
$\sim$24$\times$ parameter range. The mean Spearman across the four
\textsc{SummEval} dimensions is in $[.284, .312]$ for all four
backbones, indicating that the metric design is qualitatively
robust to embedder choice. Per-dimension profiles, however, shift
non-trivially: \texttt{nomic-embed-text} is strongest on coherence
and relevance, the alternatives on consistency and fluency. The
headline correlations in Tables~\ref{tab:correlations_summary} and
\ref{tab:compare_lgs_summary} are reported with
\texttt{nomic-embed-text} (the smallest of the higher-mean
embedders); per-dimension comparisons should therefore be read as
backbone-conditional.

\begin{table*}[t]
\centering
\caption{Embedder ablation: canonical LGS ($\lambda=0.5$) on the full SummEval set with four sentence-embedder backbones (Ollama). Reports summary-level Spearman $\rho$ per dimension and the mean across the four. Canonical embedder (used in the headline tables) marked $\star$.}
\label{tab:ablation_embedder}
\begin{tabular}{lrrrrr}
\toprule
Embedder & Coh & Con & Flu & Rel & Mean \\
\midrule
\texttt{nomic-embed-text} (137M) $\star$ & .396 & .227 & .156 & .356 & .284 \\
\texttt{mxbai-embed-large} (335M) & .304 & .351 & .263 & .266 & .296 \\
\texttt{all-minilm} (23M) & .291 & .367 & .276 & .259 & .298 \\
\texttt{bge-m3} (567M) & .335 & .360 & .251 & .299 & .312 \\
\bottomrule
\end{tabular}
\end{table*}


\subsection{Cross-embedder lead-bias verification}
\label{sec:ablation-lead-bge}
% TODO(prose): one paragraph. Key facts to land:
% - bge-m3 dev sweep peaks at lambda*=0.25, not 0.5
% - applying nomic-tuned lambda=0.5 to bge-m3 underperforms the
%   no-prior baseline on dev (.302 vs .305)
% - the qualitative shape of the lambda vs mean rho curve is
%   preserved (peak at lambda>0, monotonic decline past peak)
% - this is the strongest empirical justification for the held-out
%   methodology being necessary, not decorative

% ------------------------------------------------------------------
\section{Discussion}
\label{sec:discussion}

% TODO(prose): write each subsection.

\subsection{Cost profile}
% Quantify: with nomic-embed-text on a single CPU thread via
% Ollama, the canonical SummEval run completes in ~XX minutes
% for all 1600 (source, candidate) pairs after per-document
% caching of source-sentence embeddings. Compare to G-Eval
% (one LLM-judge call per candidate, with sampled decoding ideally
% averaged over 20 runs per the original paper) and UniEval
% (one fine-tuned T5-large forward pass per
% (article, candidate, dimension) tuple).

\subsection{What LGS does and does not claim}
% Explicit honesty paragraph aligning with the project's
% positioning. Three sub-claims:
% (i) reference-free + cheap is the contribution, not SOTA;
% (ii) per-dimension test-split lift over the no-prior baseline
%      is meaningful only on Coh and on the mean; Con/Flu/Rel
%      are essentially flat on test;
% (iii) lambda=0.5 is the dev-selected canonical with nomic;
%       practitioners using a different embedder must run their
%       own dev sweep.

\subsection{Limitations}
% (a) SummEval is CNN/DailyMail-only news; lead bias is a
%     domain-specific structural property of this corpus and
%     LGS may transfer worse to corpora without strong lead bias
%     (e.g., scientific abstracting, narrative summarisation).
% (b) G-Eval is reported as a single-run baseline (greedy
%     decoding); the original paper averages 20 samples to absorb
%     LLM-judge variance. Paired-bootstrap p-values vs G-Eval thus
%     conflate G-Eval's stochasticity with our own. Future work
%     should report G-Eval mean +- std over 3-5 runs.
% (c) The lead-bias prior is exponential; we have not tested
%     linear, reciprocal, or hard-cutoff alternatives.

% ------------------------------------------------------------------
\section{Conclusion}
\label{sec:conclusion}

% TODO(prose): 1-paragraph wrap. Summarise contribution, restate
% honest positioning vs UniEval/G-Eval, point to deployment
% scenarios that prefer LGS (no human reference, deterministic
% scoring, single-embedder cost) and to the held-out methodology
% as a transferable contribution beyond this specific metric.

% ------------------------------------------------------------------
\section*{Reproducibility}
% Code: github URL placeholder.
% Data: SummEval is publicly available.
% Embedder: nomic-embed-text via Ollama, exact tag pinned in the
% Makefile of the released artifact.
% Bootstrap seed and resample count documented in pkg/eval.

% ------------------------------------------------------------------
\bibliographystyle{IEEEtran}
% TODO(refs): create paper/refs.bib with the citations referenced
% above (fabbri2021summeval, papineni2002bleu, lin2004rouge,
% banerjee2005meteor, popovic2015chrf, zhang2020bertscore,
% zhao2019moverscore, yuan2021bartscore, amplayo2022smart,
% liu2023geval, zhong2022unieval) and replace the manual thebibliography
% block below with \bibliography{refs}.
\begin{thebibliography}{1}

\bibitem{fabbri2021summeval}
A.~R. Fabbri \emph{et al.},
``SummEval: Re-evaluating Summarization Evaluation,''
\emph{Trans. Assoc. Comput. Linguist.}, vol.~9, pp.~391--409, 2021.

\bibitem{papineni2002bleu}
K.~Papineni \emph{et al.},
``BLEU: a Method for Automatic Evaluation of Machine Translation,''
in \emph{Proc. ACL}, 2002, pp.~311--318.

\bibitem{lin2004rouge}
C.-Y. Lin,
``ROUGE: A Package for Automatic Evaluation of Summaries,''
in \emph{ACL Workshop on Text Summarization Branches Out}, 2004.

\bibitem{banerjee2005meteor}
S.~Banerjee and A.~Lavie,
``METEOR: An Automatic Metric for MT Evaluation with Improved
Correlation with Human Judgments,''
in \emph{Proc. ACL Workshop on Intrinsic and Extrinsic Evaluation
Measures for MT and/or Summarization}, 2005, pp.~65--72.

\bibitem{popovic2015chrf}
M.~Popovi\'c,
``chrF: character n-gram F-score for automatic MT evaluation,''
in \emph{Proc. WMT}, 2015, pp.~392--395.

\bibitem{zhang2020bertscore}
T.~Zhang \emph{et al.},
``BERTScore: Evaluating Text Generation with BERT,''
in \emph{Proc. ICLR}, 2020.

\bibitem{zhao2019moverscore}
W.~Zhao \emph{et al.},
``MoverScore: Text Generation Evaluating with Contextualized
Embeddings and Earth Mover Distance,''
in \emph{Proc. EMNLP-IJCNLP}, 2019, pp.~563--578.

\bibitem{yuan2021bartscore}
W.~Yuan, G.~Neubig, and P.~Liu,
``BARTScore: Evaluating Generated Text as Text Generation,''
in \emph{Proc. NeurIPS}, 2021.

\bibitem{amplayo2022smart}
R.~K. Amplayo, P.~K. Liu, and M.~Lapata,
``SMART: Sentences as Basic Units for Text Evaluation,''
\emph{arXiv preprint arXiv:2208.01030}, 2022.

\bibitem{liu2023geval}
Y.~Liu \emph{et al.},
``G-Eval: NLG Evaluation using GPT-4 with Better Human Alignment,''
in \emph{Proc. EMNLP}, 2023.

\bibitem{zhong2022unieval}
M.~Zhong \emph{et al.},
``Towards a Unified Multi-Dimensional Evaluator for Text
Generation,''
in \emph{Proc. EMNLP}, 2022, pp.~2023--2038.

\end{thebibliography}

\end{document}
